% GENERATED FILE. Edit canonical lessons, then run npm run generate:books.
% canonical-source-hash: fnv1a64:7511ff5f205064bb
% canonical-lessons: JA-C505-rikishi, JA-C505-dohyou, JA-C505-gyouji, JA-C505-yumi, JA-C505-yokeru

\chapter{Sumo and Archery}
\label{ch:ja-sumo-and-archery}
\hlchaptermodality{505}

\begin{chapteropening}
\textbf{By the end of this chapter:} \emph{I can talk about two traditional sports.}

Complete the last lesson of chapter 505: I can talk about two traditional sports.
\end{chapteropening}

\section[\texorpdfstring{\ja{りきし} (rikishi)}{りきし (rikishi)}]{\ja{りきし} (rikishi) — a sumo wrestler}

\label{lesson:JA-C505-rikishi}



\begin{quote}
Before the new one: say the Japanese for an artist, then the Japanese for a picture book.
\end{quote}

\subsection*{You'll want to know: \ja{りきし}}
\textbf{\ja{りきし}} — \emph{rikishi} — \textquotedblleft{}a sumo wrestler\textquotedblright{}.

\begin{quote}
\textbf{\ja{りきしは} \ja{おおきいです。}} — \emph{rikishi wa ōkii desu} — \textquotedblleft{}Sumo wrestlers are big.\textquotedblright{}
\end{quote}

\begin{cousinweb}
\textbf{\ja{りきし}} (\emph{rikishi}) is Sino-Japanese (kango): two Chinese characters read with their on'yomi, readings borrowed from Chinese. They are \emph{riki}, \textquotedblleft{}power\textquotedblright{} (Middle Chinese \emph{lik}) and \emph{shi}, \textquotedblleft{}gentleman\textquotedblright{} (\emph{dzriX}) (KANJIDIC2; Middle Chinese in Baxter's transcription).
\end{cousinweb}

\begin{grammarlens}[title={What you've built}]
One more word for describing what you lived and felt.
\end{grammarlens}

\subsection*{Your turn}
\begin{itemize}
\raggedright
  \item \emph{Say it:} \emph{rikishi}
  \item \emph{Say it:} \emph{rikishi}, once more
  \item \emph{Say it:} \emph{rikishi wa ōkii desu} — Sumo wrestlers are big.
  \item \emph{Recall:} say the Japanese for an artist, then the Japanese for a picture book, then say \emph{rikishi} again
\end{itemize}

\begin{tcolorbox}[breakable,colback=teal!4,colframe=teal!35!black,title={Before you move on}]
What does \ja{りきし} mean? (\textquotedblleft{}A sumo wrestler\textquotedblright{}.) Say it once more.
\end{tcolorbox}
\section[\texorpdfstring{\ja{どひょう} (dohyō)}{どひょう (dohyō)}]{\ja{どひょう} (dohyō) — a sumo ring}

\label{lesson:JA-C505-dohyou}



\begin{quote}
Before the new one: say the Japanese for a picture book, then the Japanese for a sumo wrestler.
\end{quote}

\subsection*{You'll want to know: \ja{どひょう}}
\textbf{\ja{どひょう}} — \emph{dohyō} — \textquotedblleft{}a sumo ring\textquotedblright{}.

\begin{quote}
\textbf{\ja{りきしが} \ja{どひょうに} \ja{あがります。}} — \emph{rikishi ga dohyō ni agarimasu} — \textquotedblleft{}The wrestlers step up into the ring.\textquotedblright{}
\end{quote}

\begin{cousinweb}
\textbf{\ja{どひょう}} (\emph{dohyō}) is Sino-Japanese (kango): two Chinese characters read with their on'yomi, readings borrowed from Chinese. They are \emph{do}, \textquotedblleft{}soil\textquotedblright{} (Middle Chinese \emph{thuX}) and \emph{hyō}, \textquotedblleft{}bag\textquotedblright{}, a sense the Middle Chinese rhyme books do not give for this character, so no older reading is cited for it (KANJIDIC2; Middle Chinese in Baxter's transcription).
\end{cousinweb}

\begin{grammarlens}[title={What you've built}]
One more word for describing what you lived and felt.
\end{grammarlens}

\subsection*{Your turn}
\begin{itemize}
\raggedright
  \item \emph{Say it:} \emph{dohyō}
  \item \emph{Say it:} \emph{dohyō}, once more
  \item \emph{Say it:} \emph{rikishi ga dohyō ni agarimasu} — The wrestlers step up into the ring.
  \item \emph{Recall:} say the Japanese for a picture book, then the Japanese for a sumo wrestler, then say \emph{dohyō} again
\end{itemize}

\begin{tcolorbox}[breakable,colback=teal!4,colframe=teal!35!black,title={Before you move on}]
What does \ja{どひょう} mean? (\textquotedblleft{}A sumo ring\textquotedblright{}.) Say it once more.
\end{tcolorbox}
\section[\texorpdfstring{\ja{ぎょうじ} (gyōji)}{ぎょうじ (gyōji)}]{\ja{ぎょうじ} (gyōji) — a sumo referee}

\label{lesson:JA-C505-gyouji}



\begin{quote}
Before the new one: say the Japanese for a sumo wrestler, then the Japanese for a sumo ring.
\end{quote}

\subsection*{You'll want to know: \ja{ぎょうじ}}
\textbf{\ja{ぎょうじ}} — \emph{gyōji} — \textquotedblleft{}a sumo referee\textquotedblright{}.

\begin{quote}
\textbf{\ja{ぎょうじは} \ja{どひょうの} \ja{なかに} \ja{います。}} — \emph{gyōji wa dohyō no naka ni imasu} — \textquotedblleft{}The referee stands in the sumo ring.\textquotedblright{}
\end{quote}

\begin{cousinweb}
\textbf{\ja{ぎょうじ}} (\emph{gyōji}) is Sino-Japanese (kango): two Chinese characters read with their on'yomi, readings borrowed from Chinese. They are \emph{gyō}, \textquotedblleft{}going\textquotedblright{} (Middle Chinese \emph{haeng}) and \emph{ji}, \textquotedblleft{}director\textquotedblright{} (\emph{si}) (KANJIDIC2; Middle Chinese in Baxter's transcription).
\end{cousinweb}

\begin{grammarlens}[title={What you've built}]
One more word for describing what you lived and felt.
\end{grammarlens}

\subsection*{Your turn}
\begin{itemize}
\raggedright
  \item \emph{Say it:} \emph{gyōji}
  \item \emph{Say it:} \emph{gyōji}, once more
  \item \emph{Say it:} \emph{gyōji wa dohyō no naka ni imasu} — The referee stands in the sumo ring.
  \item \emph{Recall:} say the Japanese for a sumo wrestler, then the Japanese for a sumo ring, then say \emph{gyōji} again
\end{itemize}

\begin{tcolorbox}[breakable,colback=teal!4,colframe=teal!35!black,title={Before you move on}]
What does \ja{ぎょうじ} mean? (\textquotedblleft{}A sumo referee\textquotedblright{}.) Say it once more.
\end{tcolorbox}
\section[\texorpdfstring{\ja{ゆみ} (yumi)}{ゆみ (yumi)}]{\ja{ゆみ} (yumi) — a bow (for arrows); archery}

\label{lesson:JA-C505-yumi}



\begin{quote}
Before the new one: say the Japanese for a sumo ring, then the Japanese for a sumo referee.
\end{quote}

\subsection*{You'll want to know: \ja{ゆみ}}
\textbf{\ja{ゆみ}} — \emph{yumi} — \textquotedblleft{}a bow (for arrows); archery\textquotedblright{}.

\begin{quote}
\textbf{\ja{おとうとは} \ja{ゆみを} \ja{ならっています。}} — \emph{otōto wa yumi o naratteimasu} — \textquotedblleft{}My younger brother is learning archery.\textquotedblright{}
\end{quote}

\begin{cousinweb}
\textbf{\ja{ゆみ}} (\emph{yumi}) is a native Japanese word (wago), inherited rather than borrowed from Chinese; its kanji was attached later as a kun'yomi, and its older origin is not settled (JMdict, KANJIDIC2).
\end{cousinweb}

\begin{grammarlens}[title={What you've built}]
One more word for describing what you lived and felt.
\end{grammarlens}

\subsection*{Your turn}
\begin{itemize}
\raggedright
  \item \emph{Say it:} \emph{yumi}
  \item \emph{Say it:} \emph{yumi}, once more
  \item \emph{Say it:} \emph{otōto wa yumi o naratteimasu} — My younger brother is learning archery.
  \item \emph{Recall:} say the Japanese for a sumo ring, then the Japanese for a sumo referee, then say \emph{yumi} again
\end{itemize}

\begin{tcolorbox}[breakable,colback=teal!4,colframe=teal!35!black,title={Before you move on}]
What does \ja{ゆみ} mean? (\textquotedblleft{}A bow (for arrows); archery\textquotedblright{}.) Say it once more.
\end{tcolorbox}
\section[\texorpdfstring{\ja{よける} (yokeru)}{よける (yokeru)}]{\ja{よける} (yokeru) — to dodge; to keep out of the way of}

\label{lesson:JA-C505-yokeru}



\begin{quote}
Before the new one: say the Japanese for a sumo referee, then the Japanese for a bow (for arrows).
\end{quote}

\subsection*{You'll want to know: \ja{よける}}
\textbf{\ja{よける}} — \emph{yokeru} — \textquotedblleft{}to dodge; to keep out of the way of\textquotedblright{}.

\begin{quote}
\textbf{\ja{いしを} \ja{よけて} \ja{あるきます。}} — \emph{ishi o yokete arukimasu} — \textquotedblleft{}I walk around the stones.\textquotedblright{}
\end{quote}

Polite \textbf{\ja{よけます}}, \emph{yokemasu}; an ichidan verb. Past (polite): \textbf{\ja{ねこが} \ja{くるまを} \ja{よけました。}}, \emph{neko ga kuruma o yokemashita}.

\begin{cousinweb}
\textbf{\ja{よける}} (\emph{yokeru}) is a native Japanese word (wago), inherited rather than borrowed from Chinese; its kanji was attached later as a kun'yomi, and its older origin is not settled (JMdict, KANJIDIC2).
\end{cousinweb}

\begin{grammarlens}[title={What you've built}]
One more word for describing what you lived and felt.
\end{grammarlens}

\subsection*{Your turn}
\begin{itemize}
\raggedright
  \item \emph{Say it:} \emph{yokeru}
  \item \emph{Say it:} \emph{yokeru}, once more
  \item \emph{Say it:} \emph{ishi o yokete arukimasu} — I walk around the stones.
  \item \emph{Recall:} say the Japanese for a sumo referee, then the Japanese for a bow (for arrows), then say \emph{yokeru} again
\end{itemize}

\begin{tcolorbox}[breakable,colback=teal!4,colframe=teal!35!black,title={Before you move on}]
What does \ja{よける} mean? (\textquotedblleft{}To dodge; to keep out of the way of\textquotedblright{}.) Say it once more.
\end{tcolorbox}
